% GENERATED FILE. Edit canonical lessons, then run npm run generate:books.
% canonical-source-hash: fnv1a64:9404559bccf37955
% canonical-lessons: TE-C471-paryavaranam, TE-C471-paccadanam, TE-C471-naruku, TE-C471-nastaparacu, TE-C471-samraksincu

\chapter{Caring for the Environment}
\label{ch:te-caring-for-the-environment}
\hlchaptermodality{471}

\begin{chapteropening}
\textbf{By the end of this chapter:} \emph{I can say the environment, greenery, to chop down, to damage and to conserve.}

Complete the last lesson of chapter 471: I can say the environment, greenery, to chop down, to damage and to conserve.
\end{chapteropening}

\section[\texorpdfstring{\te{పర్యావరణం} (paryāvaraṇaṁ)}{పర్యావరణం (paryāvaraṇaṁ)}]{\te{పర్యావరణం} (paryāvaraṇaṁ) — the environment}

\label{lesson:TE-C471-paryavaranam}



\begin{quote}
Before the new one: say the Telugu for to pollute, then the Telugu for pure.
\end{quote}

\subsection*{You'll want to know: \te{పర్యావరణం}}
\textbf{\te{పర్యావరణం}} — \emph{paryāvaraṇaṁ} — \textquotedblleft{}the environment\textquotedblright{}.

\begin{quote}
\te{పర్యావరణం} \te{కోసం} \te{ప్లాస్టిక్} \te{వాడకం} \te{తగ్గించాలి}.
\end{quote}

\begin{quote}
\emph{paryāvaraṇaṁ kōsaṁ plāsṭik vāḍakaṁ taggiñcāli.}
\end{quote}

\begin{quote}
For the sake of the environment we should use less plastic.
\end{quote}

\te{పర్యావరణ}, environmental, before a noun: \te{పర్యావరణ} \te{దినోత్సవం}, Environment Day.

\begin{cousinweb}
\textbf{\te{పర్యావరణం}} (\emph{paryāvaraṇaṁ}) is a modern Telugu word built from Sanskrit \emph{pari}, \textquotedblleft{}around\textquotedblright{}, and \emph{āvaraṇa}, \textquotedblleft{}a covering, an enclosure\textquotedblright{} (Monier-Williams): what surrounds and covers us.
\end{cousinweb}

\begin{grammarlens}[title={What you've built}]
One more word for getting around.
\end{grammarlens}

\subsection*{Your turn}
\begin{itemize}
\raggedright
  \item \emph{Say it:} \emph{paryāvaraṇaṁ}
  \item \emph{Say it:} \emph{paryāvaraṇaṁ}, once more
  \item \emph{Say it:} \emph{paryāvaraṇaṁ kōsaṁ plāsṭik vāḍakaṁ taggiñcāli.} — For the sake of the environment we should use less plastic.
  \item \emph{Recall:} say the Telugu for to pollute, then the Telugu for pure, then say \emph{paryāvaraṇaṁ} again
\end{itemize}

\begin{tcolorbox}[breakable,colback=teal!4,colframe=teal!35!black,title={Before you move on}]
What does \te{పర్యావరణం} mean? (\textquotedblleft{}The environment\textquotedblright{}.) Say it once more.
\end{tcolorbox}
\section[\texorpdfstring{\te{పచ్చదనం} (paccadanaṁ)}{పచ్చదనం (paccadanaṁ)}]{\te{పచ్చదనం} (paccadanaṁ) — greenery}

\label{lesson:TE-C471-paccadanam}



\begin{quote}
Before the new one: say the Telugu for pure, then the Telugu for the environment.
\end{quote}

\subsection*{You'll want to know: \te{పచ్చదనం}}
\textbf{\te{పచ్చదనం}} — \emph{paccadanaṁ} — \textquotedblleft{}greenery\textquotedblright{}.

\begin{quote}
\te{వర్షాల} \te{తరువాత} \te{పొలాల్లో} \te{పచ్చదనం} \te{కనిపిస్తోంది}.
\end{quote}

\begin{quote}
\emph{varṣāla taruvāta polāllō paccadanaṁ kanipistōndi.}
\end{quote}

\begin{quote}
After the rains greenery is showing in the fields.
\end{quote}

\te{పచ్చ}, green, + -\te{దనం}, the ending of an abstract noun: greenness.

\begin{cousinweb}
\textbf{\te{పచ్చదనం}} (\emph{paccadanaṁ}) is native Dravidian: \te{పచ్చ}, \textquotedblleft{}green\textquotedblright{} (DEDR 3821), with the abstract ending -\te{దనం}.
\end{cousinweb}

\begin{grammarlens}[title={What you've built}]
One more word for getting around.
\end{grammarlens}

\subsection*{Your turn}
\begin{itemize}
\raggedright
  \item \emph{Say it:} \emph{paccadanaṁ}
  \item \emph{Say it:} \emph{paccadanaṁ}, once more
  \item \emph{Say it:} \emph{varṣāla taruvāta polāllō paccadanaṁ kanipistōndi.} — After the rains greenery is showing in the fields.
  \item \emph{Recall:} say the Telugu for pure, then the Telugu for the environment, then say \emph{paccadanaṁ} again
\end{itemize}

\begin{tcolorbox}[breakable,colback=teal!4,colframe=teal!35!black,title={Before you move on}]
What does \te{పచ్చదనం} mean? (\textquotedblleft{}Greenery\textquotedblright{}.) Say it once more.
\end{tcolorbox}
\section[\texorpdfstring{\te{నరుకు} (naruku)}{నరుకు (naruku)}]{\te{నరుకు} (naruku) — to chop down; to fell}

\label{lesson:TE-C471-naruku}



\begin{quote}
Before the new one: say the Telugu for the environment, then the Telugu for greenery.
\end{quote}

\subsection*{You'll want to know: \te{నరుకు}}
\textbf{\te{నరుకు}} — \emph{naruku} — \textquotedblleft{}to chop down; to fell\textquotedblright{}.

\begin{quote}
\te{రోడ్డు} \te{కోసం} \te{చాలా} \te{చెట్లు} \te{నరికారు}.
\end{quote}

\begin{quote}
\emph{rōḍḍu kōsaṁ cālā ceṭlu narikāru.}
\end{quote}

\begin{quote}
Many trees were cut down for the road.
\end{quote}

Past (I …): \te{నరికాను}, \emph{narikānu}.

Said of trees and of wood chopped with an axe.

\begin{cousinweb}
\textbf{\te{నరుకు}} (\emph{naruku}) is native Dravidian: DEDR 3625 lists Telugu \emph{naṟuku}, \textquotedblleft{}to cut, to hew, to fell\textquotedblright{}, with Tamil \emph{naṟukku}, \textquotedblleft{}to chop\textquotedblright{}.
\end{cousinweb}

\begin{grammarlens}[title={What you've built}]
One more word for getting around.
\end{grammarlens}

\subsection*{Your turn}
\begin{itemize}
\raggedright
  \item \emph{Say it:} \emph{naruku}
  \item \emph{Say it:} \emph{naruku}, once more
  \item \emph{Say it:} \emph{rōḍḍu kōsaṁ cālā ceṭlu narikāru.} — Many trees were cut down for the road.
  \item \emph{Recall:} say the Telugu for the environment, then the Telugu for greenery, then say \emph{naruku} again
\end{itemize}

\begin{tcolorbox}[breakable,colback=teal!4,colframe=teal!35!black,title={Before you move on}]
What does \te{నరుకు} mean? (\textquotedblleft{}To chop down; to fell\textquotedblright{}.) Say it once more.
\end{tcolorbox}
\section[\texorpdfstring{\te{నష్టపరచు} (naṣṭaparacu)}{నష్టపరచు (naṣṭaparacu)}]{\te{నష్టపరచు} (naṣṭaparacu) — to damage; to harm}

\label{lesson:TE-C471-nastaparacu}



\begin{quote}
Before the new one: say the Telugu for greenery, then the Telugu for to chop down.
\end{quote}

\subsection*{You'll want to know: \te{నష్టపరచు}}
\textbf{\te{నష్టపరచు}} — \emph{naṣṭaparacu} — \textquotedblleft{}to damage; to harm\textquotedblright{}.

\begin{quote}
\te{వరదలు} \te{పంటలను} \te{నష్టపరిచాయి}.
\end{quote}

\begin{quote}
\emph{varadalu paṇṭalanu naṣṭaparicāyi.}
\end{quote}

\begin{quote}
The floods damaged the crops.
\end{quote}

Past (I …): \te{నష్టపరిచాను}, \emph{naṣṭaparicānu}.

\te{నష్టం}, a loss, + -\te{పరచు}, to cause: to cause a loss. The same -\te{పరచు} makes other verbs from nouns.

\begin{cousinweb}
\textbf{\te{నష్టపరచు}} (\emph{naṣṭaparacu}) is a Telugu verb built from \te{నష్టం}, \textquotedblleft{}a loss\textquotedblright{}, borrowed from Sanskrit \emph{naṣṭa}, \textquotedblleft{}lost, destroyed\textquotedblright{} (Monier-Williams), and the native -\te{పరచు}, the causative of \te{పడు}, which DEDR 3853 lists as \emph{paṟacu}, \textquotedblleft{}to cause to feel or suffer\textquotedblright{}.
\end{cousinweb}

\begin{grammarlens}[title={What you've built}]
One more word for getting around.
\end{grammarlens}

\subsection*{Your turn}
\begin{itemize}
\raggedright
  \item \emph{Say it:} \emph{naṣṭaparacu}
  \item \emph{Say it:} \emph{naṣṭaparacu}, once more
  \item \emph{Say it:} \emph{varadalu paṇṭalanu naṣṭaparicāyi.} — The floods damaged the crops.
  \item \emph{Recall:} say the Telugu for greenery, then the Telugu for to chop down, then say \emph{naṣṭaparacu} again
\end{itemize}

\begin{tcolorbox}[breakable,colback=teal!4,colframe=teal!35!black,title={Before you move on}]
What does \te{నష్టపరచు} mean? (\textquotedblleft{}To damage; to harm\textquotedblright{}.) Say it once more.
\end{tcolorbox}
\section[\texorpdfstring{\te{సంరక్షించు} (saṁrakṣiñcu)}{సంరక్షించు (saṁrakṣiñcu)}]{\te{సంరక్షించు} (saṁrakṣiñcu) — to conserve; to protect}

\label{lesson:TE-C471-samraksincu}



\begin{quote}
Before the new one: say the Telugu for to chop down, then the Telugu for to damage.
\end{quote}

\subsection*{You'll want to know: \te{సంరక్షించు}}
\textbf{\te{సంరక్షించు}} — \emph{saṁrakṣiñcu} — \textquotedblleft{}to conserve; to protect\textquotedblright{}.

\begin{quote}
\te{గ్రామస్తులు} \te{ఆ} \te{అడవిని} \te{సంరక్షిస్తున్నారు}.
\end{quote}

\begin{quote}
\emph{grāmastulu ā aḍavini saṁrakṣistunnāru.}
\end{quote}

\begin{quote}
The villagers are protecting that forest.
\end{quote}

Past (I …): \te{సంరక్షించాను}, \emph{saṁrakṣiñcānu}.

A formal word; the noun is \te{సంరక్షణ}, conservation.

\begin{cousinweb}
\textbf{\te{సంరక్షించు}} (\emph{saṁrakṣiñcu}) is built from Sanskrit \emph{saṁ-rakṣ}, \textquotedblleft{}to guard, to preserve\textquotedblright{}, whose noun is \emph{saṁrakṣaṇa}, \textquotedblleft{}preservation\textquotedblright{} (Monier-Williams), with the Telugu verb ending -\te{ఇంచు}.
\end{cousinweb}

\begin{grammarlens}[title={What you've built}]
One more word for getting around.
\end{grammarlens}

\subsection*{Your turn}
\begin{itemize}
\raggedright
  \item \emph{Say it:} \emph{saṁrakṣiñcu}
  \item \emph{Say it:} \emph{saṁrakṣiñcu}, once more
  \item \emph{Say it:} \emph{grāmastulu ā aḍavini saṁrakṣistunnāru.} — The villagers are protecting that forest.
  \item \emph{Recall:} say the Telugu for to chop down, then the Telugu for to damage, then say \emph{saṁrakṣiñcu} again
\end{itemize}

\begin{tcolorbox}[breakable,colback=teal!4,colframe=teal!35!black,title={Before you move on}]
What does \te{సంరక్షించు} mean? (\textquotedblleft{}To conserve; to protect\textquotedblright{}.) Say it once more.
\end{tcolorbox}
